\section{主要结果}

\subsection{相同的设计策略能否在有限角度域上维持目标算子形状？}

锚定测试是用于二阶傅里叶微分的自由形态超表面。这个目标是有意苛刻的。成功的器件必须抑制低空间频率中心，向更大的横向波矢量增加其响应，并在有限的入射角区间而非单一名义点上这样做。在项目数据中，二阶排序逐波长评估，根据结合中心抑制、目标形状一致性和边缘响应的分数选择顶级候选。结果分布已经揭示有意义的设计景观：在采样波长上，有效候选池包含低质量结构的长尾和得分高得多的几何的小得多的子集，这是稀疏但非空可行集的特征，而非人为容易的设计问题。

最强的二阶候选在物理上很有趣有两个原因。首先，它们的验证响应保留了与二阶导数传递函数相关的预期偶对称性和向外增长。其次，这种行为不限于档案中的一个孤立波长。排序摘要和每波长顶级候选图表明，相同的响应中心设计逻辑在有限工作窗口上保持有意义。因此，目标更好地被理解为角度和波长上的结构化算子规范，而非单点曲线拟合练习。

\subsection{替代预测和验证物理是否在同一响应对象上一致？}

对于任何数据驱动的逆向设计程序，一个关键问题是习得的近似和物理求解器是否以足够相似的方式评估相同的光学对象。这里该对象是响应场本身。项目使这种区别明确：替代模型首先估计目标条件响应质量，但最终图和分数在RCWA评估后重新计算。这种分离是必要的，因为响应场设计可能以微妙的方式失败。候选可能在中心附近保持目标曲率，同时扭曲肩部；它在简化指标下可能看起来有希望，同时在所需角度支持外泄漏功率；或者它可能仅在归一化后与目标对齐，掩盖对器件操作仍相关的幅度缺陷。

验证的二阶结果显示为什么完整响应场比较比单一数字报告更有信息量。仓库中使用的分数分量编码物理上重要的特征：零空间频率附近的低透射、期望二阶包络的恢复、以及设计孔径上有用边缘响应的保持。具有相似名义总分数的候选仍可在这些贡献如何平衡上显著不同。这正是响应空间视角变得有用的地方。通过比较验证场而非局部性能数字，可以区分仅在一个条件附近类似于目标的设计与在预期角度跨度上维持正确算子几何的设计。

同样的原理解释为什么生成后重排序不能省略。如果替代一致性足够，物理和名义排序将与噪声重合。实际上它们不重合。项目保留单独的验证和基于RCWA的重排序阶段，正是因为逆图足够丰富，小的几何差异可以在验证算子场中产生重要差异。在这种背景下，物理一致性不是哲学偏好。它是使最终设计与散射物理而非习得近似相关联的操作规则。

\subsection{响应级一致性是否转化为实际的算子行为？}

算子超表面最终由它们对光场和图像的作用来判断，而不仅由它们的分数曲线看起来多平滑来判断。因此，图像域测试作为次要但重要的验证层包含在项目中。二阶目标特别适合此目的，因为它对输入图像的预期作用是直观的：抑制宽泛的低频内容，同时增强曲率敏感特征和边缘状结构。当验证响应用于图像处理演示时，结果输出恢复特征性的二阶导数行为，而非任意的对比度增强伪影。

从响应级一致性到算子级行为的这种转化闭合了本节科学论证。强大的逆向设计结果不仅是针对内部指标得分高的结构。它是RCWA验证响应具有正确算子形状且对代表性输入的作用与该形状一致的结构。因此，这个大NA、高透过率、宽带代表性例子以精确方式支持本文中心论点：以目标响应场为条件、采样多个候选、并通过物理验证选择它们，足以恢复实际有意义的傅里叶算子超表面，而无需将问题简化为手工专门化的几何族。
